\documentclass[11pt]{article}

    \usepackage[T2A]{fontenc}
\usepackage[utf8]{inputenc}
\usepackage[english, russian]{babel}

\usepackage[breakable]{tcolorbox}
    \usepackage{parskip} % Stop auto-indenting (to mimic markdown behaviour)
    
    \usepackage{iftex}
    \ifPDFTeX
    	\usepackage[T1]{fontenc}
    	\usepackage{mathpazo}
    \else
    	\usepackage{fontspec}
    \fi

    % Basic figure setup, for now with no caption control since it's done
    % automatically by Pandoc (which extracts ![](path) syntax from Markdown).
    \usepackage{graphicx}
    % Maintain compatibility with old templates. Remove in nbconvert 6.0
    \let\Oldincludegraphics\includegraphics
    % Ensure that by default, figures have no caption (until we provide a
    % proper Figure object with a Caption API and a way to capture that
    % in the conversion process - todo).
    \usepackage{caption}
    \DeclareCaptionFormat{nocaption}{}
    \captionsetup{format=nocaption,aboveskip=0pt,belowskip=0pt}

    \usepackage{float}
    \floatplacement{figure}{H} % forces figures to be placed at the correct location
    \usepackage{xcolor} % Allow colors to be defined
    \usepackage{enumerate} % Needed for markdown enumerations to work
    \usepackage{geometry} % Used to adjust the document margins
    \usepackage{amsmath} % Equations
    \usepackage{amssymb} % Equations
    \usepackage{textcomp} % defines textquotesingle
    % Hack from http://tex.stackexchange.com/a/47451/13684:
    \AtBeginDocument{%
        \def\PYZsq{\textquotesingle}% Upright quotes in Pygmentized code
    }
    \usepackage{upquote} % Upright quotes for verbatim code
    \usepackage{eurosym} % defines \euro
    \usepackage[mathletters]{ucs} % Extended unicode (utf-8) support
    \usepackage{fancyvrb} % verbatim replacement that allows latex
    \usepackage{grffile} % extends the file name processing of package graphics 
                         % to support a larger range
    \makeatletter % fix for old versions of grffile with XeLaTeX
    \@ifpackagelater{grffile}{2019/11/01}
    {
      % Do nothing on new versions
    }
    {
      \def\Gread@@xetex#1{%
        \IfFileExists{"\Gin@base".bb}%
        {\Gread@eps{\Gin@base.bb}}%
        {\Gread@@xetex@aux#1}%
      }
    }
    \makeatother
    \usepackage[Export]{adjustbox} % Used to constrain images to a maximum size
    \adjustboxset{max size={0.9\linewidth}{0.9\paperheight}}

    % The hyperref package gives us a pdf with properly built
    % internal navigation ('pdf bookmarks' for the table of contents,
    % internal cross-reference links, web links for URLs, etc.)
    \usepackage{hyperref}
    % The default LaTeX title has an obnoxious amount of whitespace. By default,
    % titling removes some of it. It also provides customization options.
    \usepackage{titling}
    \usepackage{longtable} % longtable support required by pandoc >1.10
    \usepackage{booktabs}  % table support for pandoc > 1.12.2
    \usepackage[inline]{enumitem} % IRkernel/repr support (it uses the enumerate* environment)
    \usepackage[normalem]{ulem} % ulem is needed to support strikethroughs (\sout)
                                % normalem makes italics be italics, not underlines
    \usepackage{mathrsfs}
    
\usepackage{wrapfig}
\usepackage[rightcaption]{sidecap}
\providecommand{\keywords}[1]{\textbf{\textit{Keywords:}} #1}

\author{A. Yu. Drozdov}

    
    % Colors for the hyperref package
    \definecolor{urlcolor}{rgb}{0,.145,.698}
    \definecolor{linkcolor}{rgb}{.71,0.21,0.01}
    \definecolor{citecolor}{rgb}{.12,.54,.11}

    % ANSI colors
    \definecolor{ansi-black}{HTML}{3E424D}
    \definecolor{ansi-black-intense}{HTML}{282C36}
    \definecolor{ansi-red}{HTML}{E75C58}
    \definecolor{ansi-red-intense}{HTML}{B22B31}
    \definecolor{ansi-green}{HTML}{00A250}
    \definecolor{ansi-green-intense}{HTML}{007427}
    \definecolor{ansi-yellow}{HTML}{DDB62B}
    \definecolor{ansi-yellow-intense}{HTML}{B27D12}
    \definecolor{ansi-blue}{HTML}{208FFB}
    \definecolor{ansi-blue-intense}{HTML}{0065CA}
    \definecolor{ansi-magenta}{HTML}{D160C4}
    \definecolor{ansi-magenta-intense}{HTML}{A03196}
    \definecolor{ansi-cyan}{HTML}{60C6C8}
    \definecolor{ansi-cyan-intense}{HTML}{258F8F}
    \definecolor{ansi-white}{HTML}{C5C1B4}
    \definecolor{ansi-white-intense}{HTML}{A1A6B2}
    \definecolor{ansi-default-inverse-fg}{HTML}{FFFFFF}
    \definecolor{ansi-default-inverse-bg}{HTML}{000000}

    % common color for the border for error outputs.
    \definecolor{outerrorbackground}{HTML}{FFDFDF}

    % commands and environments needed by pandoc snippets
    % extracted from the output of `pandoc -s`
    \providecommand{\tightlist}{%
      \setlength{\itemsep}{0pt}\setlength{\parskip}{0pt}}
    \DefineVerbatimEnvironment{Highlighting}{Verbatim}{commandchars=\\\{\}}
    % Add ',fontsize=\small' for more characters per line
    \newenvironment{Shaded}{}{}
    \newcommand{\KeywordTok}[1]{\textcolor[rgb]{0.00,0.44,0.13}{\textbf{{#1}}}}
    \newcommand{\DataTypeTok}[1]{\textcolor[rgb]{0.56,0.13,0.00}{{#1}}}
    \newcommand{\DecValTok}[1]{\textcolor[rgb]{0.25,0.63,0.44}{{#1}}}
    \newcommand{\BaseNTok}[1]{\textcolor[rgb]{0.25,0.63,0.44}{{#1}}}
    \newcommand{\FloatTok}[1]{\textcolor[rgb]{0.25,0.63,0.44}{{#1}}}
    \newcommand{\CharTok}[1]{\textcolor[rgb]{0.25,0.44,0.63}{{#1}}}
    \newcommand{\StringTok}[1]{\textcolor[rgb]{0.25,0.44,0.63}{{#1}}}
    \newcommand{\CommentTok}[1]{\textcolor[rgb]{0.38,0.63,0.69}{\textit{{#1}}}}
    \newcommand{\OtherTok}[1]{\textcolor[rgb]{0.00,0.44,0.13}{{#1}}}
    \newcommand{\AlertTok}[1]{\textcolor[rgb]{1.00,0.00,0.00}{\textbf{{#1}}}}
    \newcommand{\FunctionTok}[1]{\textcolor[rgb]{0.02,0.16,0.49}{{#1}}}
    \newcommand{\RegionMarkerTok}[1]{{#1}}
    \newcommand{\ErrorTok}[1]{\textcolor[rgb]{1.00,0.00,0.00}{\textbf{{#1}}}}
    \newcommand{\NormalTok}[1]{{#1}}
    
    % Additional commands for more recent versions of Pandoc
    \newcommand{\ConstantTok}[1]{\textcolor[rgb]{0.53,0.00,0.00}{{#1}}}
    \newcommand{\SpecialCharTok}[1]{\textcolor[rgb]{0.25,0.44,0.63}{{#1}}}
    \newcommand{\VerbatimStringTok}[1]{\textcolor[rgb]{0.25,0.44,0.63}{{#1}}}
    \newcommand{\SpecialStringTok}[1]{\textcolor[rgb]{0.73,0.40,0.53}{{#1}}}
    \newcommand{\ImportTok}[1]{{#1}}
    \newcommand{\DocumentationTok}[1]{\textcolor[rgb]{0.73,0.13,0.13}{\textit{{#1}}}}
    \newcommand{\AnnotationTok}[1]{\textcolor[rgb]{0.38,0.63,0.69}{\textbf{\textit{{#1}}}}}
    \newcommand{\CommentVarTok}[1]{\textcolor[rgb]{0.38,0.63,0.69}{\textbf{\textit{{#1}}}}}
    \newcommand{\VariableTok}[1]{\textcolor[rgb]{0.10,0.09,0.49}{{#1}}}
    \newcommand{\ControlFlowTok}[1]{\textcolor[rgb]{0.00,0.44,0.13}{\textbf{{#1}}}}
    \newcommand{\OperatorTok}[1]{\textcolor[rgb]{0.40,0.40,0.40}{{#1}}}
    \newcommand{\BuiltInTok}[1]{{#1}}
    \newcommand{\ExtensionTok}[1]{{#1}}
    \newcommand{\PreprocessorTok}[1]{\textcolor[rgb]{0.74,0.48,0.00}{{#1}}}
    \newcommand{\AttributeTok}[1]{\textcolor[rgb]{0.49,0.56,0.16}{{#1}}}
    \newcommand{\InformationTok}[1]{\textcolor[rgb]{0.38,0.63,0.69}{\textbf{\textit{{#1}}}}}
    \newcommand{\WarningTok}[1]{\textcolor[rgb]{0.38,0.63,0.69}{\textbf{\textit{{#1}}}}}
    
    
    % Define a nice break command that doesn't care if a line doesn't already
    % exist.
    \def\br{\hspace*{\fill} \\* }
    % Math Jax compatibility definitions
    \def\gt{>}
    \def\lt{<}
    \let\Oldtex\TeX
    \let\Oldlatex\LaTeX
    \renewcommand{\TeX}{\textrm{\Oldtex}}
    \renewcommand{\LaTeX}{\textrm{\Oldlatex}}
    % Document parameters
    % Document title
    \title{03.do\_lagging\_potentials\_equations\_correspond\_to\_Lorentz\_calibration}
    
    
    
    
    
% Pygments definitions
\makeatletter
\def\PY@reset{\let\PY@it=\relax \let\PY@bf=\relax%
    \let\PY@ul=\relax \let\PY@tc=\relax%
    \let\PY@bc=\relax \let\PY@ff=\relax}
\def\PY@tok#1{\csname PY@tok@#1\endcsname}
\def\PY@toks#1+{\ifx\relax#1\empty\else%
    \PY@tok{#1}\expandafter\PY@toks\fi}
\def\PY@do#1{\PY@bc{\PY@tc{\PY@ul{%
    \PY@it{\PY@bf{\PY@ff{#1}}}}}}}
\def\PY#1#2{\PY@reset\PY@toks#1+\relax+\PY@do{#2}}

\@namedef{PY@tok@w}{\def\PY@tc##1{\textcolor[rgb]{0.73,0.73,0.73}{##1}}}
\@namedef{PY@tok@c}{\let\PY@it=\textit\def\PY@tc##1{\textcolor[rgb]{0.25,0.50,0.50}{##1}}}
\@namedef{PY@tok@cp}{\def\PY@tc##1{\textcolor[rgb]{0.74,0.48,0.00}{##1}}}
\@namedef{PY@tok@k}{\let\PY@bf=\textbf\def\PY@tc##1{\textcolor[rgb]{0.00,0.50,0.00}{##1}}}
\@namedef{PY@tok@kp}{\def\PY@tc##1{\textcolor[rgb]{0.00,0.50,0.00}{##1}}}
\@namedef{PY@tok@kt}{\def\PY@tc##1{\textcolor[rgb]{0.69,0.00,0.25}{##1}}}
\@namedef{PY@tok@o}{\def\PY@tc##1{\textcolor[rgb]{0.40,0.40,0.40}{##1}}}
\@namedef{PY@tok@ow}{\let\PY@bf=\textbf\def\PY@tc##1{\textcolor[rgb]{0.67,0.13,1.00}{##1}}}
\@namedef{PY@tok@nb}{\def\PY@tc##1{\textcolor[rgb]{0.00,0.50,0.00}{##1}}}
\@namedef{PY@tok@nf}{\def\PY@tc##1{\textcolor[rgb]{0.00,0.00,1.00}{##1}}}
\@namedef{PY@tok@nc}{\let\PY@bf=\textbf\def\PY@tc##1{\textcolor[rgb]{0.00,0.00,1.00}{##1}}}
\@namedef{PY@tok@nn}{\let\PY@bf=\textbf\def\PY@tc##1{\textcolor[rgb]{0.00,0.00,1.00}{##1}}}
\@namedef{PY@tok@ne}{\let\PY@bf=\textbf\def\PY@tc##1{\textcolor[rgb]{0.82,0.25,0.23}{##1}}}
\@namedef{PY@tok@nv}{\def\PY@tc##1{\textcolor[rgb]{0.10,0.09,0.49}{##1}}}
\@namedef{PY@tok@no}{\def\PY@tc##1{\textcolor[rgb]{0.53,0.00,0.00}{##1}}}
\@namedef{PY@tok@nl}{\def\PY@tc##1{\textcolor[rgb]{0.63,0.63,0.00}{##1}}}
\@namedef{PY@tok@ni}{\let\PY@bf=\textbf\def\PY@tc##1{\textcolor[rgb]{0.60,0.60,0.60}{##1}}}
\@namedef{PY@tok@na}{\def\PY@tc##1{\textcolor[rgb]{0.49,0.56,0.16}{##1}}}
\@namedef{PY@tok@nt}{\let\PY@bf=\textbf\def\PY@tc##1{\textcolor[rgb]{0.00,0.50,0.00}{##1}}}
\@namedef{PY@tok@nd}{\def\PY@tc##1{\textcolor[rgb]{0.67,0.13,1.00}{##1}}}
\@namedef{PY@tok@s}{\def\PY@tc##1{\textcolor[rgb]{0.73,0.13,0.13}{##1}}}
\@namedef{PY@tok@sd}{\let\PY@it=\textit\def\PY@tc##1{\textcolor[rgb]{0.73,0.13,0.13}{##1}}}
\@namedef{PY@tok@si}{\let\PY@bf=\textbf\def\PY@tc##1{\textcolor[rgb]{0.73,0.40,0.53}{##1}}}
\@namedef{PY@tok@se}{\let\PY@bf=\textbf\def\PY@tc##1{\textcolor[rgb]{0.73,0.40,0.13}{##1}}}
\@namedef{PY@tok@sr}{\def\PY@tc##1{\textcolor[rgb]{0.73,0.40,0.53}{##1}}}
\@namedef{PY@tok@ss}{\def\PY@tc##1{\textcolor[rgb]{0.10,0.09,0.49}{##1}}}
\@namedef{PY@tok@sx}{\def\PY@tc##1{\textcolor[rgb]{0.00,0.50,0.00}{##1}}}
\@namedef{PY@tok@m}{\def\PY@tc##1{\textcolor[rgb]{0.40,0.40,0.40}{##1}}}
\@namedef{PY@tok@gh}{\let\PY@bf=\textbf\def\PY@tc##1{\textcolor[rgb]{0.00,0.00,0.50}{##1}}}
\@namedef{PY@tok@gu}{\let\PY@bf=\textbf\def\PY@tc##1{\textcolor[rgb]{0.50,0.00,0.50}{##1}}}
\@namedef{PY@tok@gd}{\def\PY@tc##1{\textcolor[rgb]{0.63,0.00,0.00}{##1}}}
\@namedef{PY@tok@gi}{\def\PY@tc##1{\textcolor[rgb]{0.00,0.63,0.00}{##1}}}
\@namedef{PY@tok@gr}{\def\PY@tc##1{\textcolor[rgb]{1.00,0.00,0.00}{##1}}}
\@namedef{PY@tok@ge}{\let\PY@it=\textit}
\@namedef{PY@tok@gs}{\let\PY@bf=\textbf}
\@namedef{PY@tok@gp}{\let\PY@bf=\textbf\def\PY@tc##1{\textcolor[rgb]{0.00,0.00,0.50}{##1}}}
\@namedef{PY@tok@go}{\def\PY@tc##1{\textcolor[rgb]{0.53,0.53,0.53}{##1}}}
\@namedef{PY@tok@gt}{\def\PY@tc##1{\textcolor[rgb]{0.00,0.27,0.87}{##1}}}
\@namedef{PY@tok@err}{\def\PY@bc##1{{\setlength{\fboxsep}{\string -\fboxrule}\fcolorbox[rgb]{1.00,0.00,0.00}{1,1,1}{\strut ##1}}}}
\@namedef{PY@tok@kc}{\let\PY@bf=\textbf\def\PY@tc##1{\textcolor[rgb]{0.00,0.50,0.00}{##1}}}
\@namedef{PY@tok@kd}{\let\PY@bf=\textbf\def\PY@tc##1{\textcolor[rgb]{0.00,0.50,0.00}{##1}}}
\@namedef{PY@tok@kn}{\let\PY@bf=\textbf\def\PY@tc##1{\textcolor[rgb]{0.00,0.50,0.00}{##1}}}
\@namedef{PY@tok@kr}{\let\PY@bf=\textbf\def\PY@tc##1{\textcolor[rgb]{0.00,0.50,0.00}{##1}}}
\@namedef{PY@tok@bp}{\def\PY@tc##1{\textcolor[rgb]{0.00,0.50,0.00}{##1}}}
\@namedef{PY@tok@fm}{\def\PY@tc##1{\textcolor[rgb]{0.00,0.00,1.00}{##1}}}
\@namedef{PY@tok@vc}{\def\PY@tc##1{\textcolor[rgb]{0.10,0.09,0.49}{##1}}}
\@namedef{PY@tok@vg}{\def\PY@tc##1{\textcolor[rgb]{0.10,0.09,0.49}{##1}}}
\@namedef{PY@tok@vi}{\def\PY@tc##1{\textcolor[rgb]{0.10,0.09,0.49}{##1}}}
\@namedef{PY@tok@vm}{\def\PY@tc##1{\textcolor[rgb]{0.10,0.09,0.49}{##1}}}
\@namedef{PY@tok@sa}{\def\PY@tc##1{\textcolor[rgb]{0.73,0.13,0.13}{##1}}}
\@namedef{PY@tok@sb}{\def\PY@tc##1{\textcolor[rgb]{0.73,0.13,0.13}{##1}}}
\@namedef{PY@tok@sc}{\def\PY@tc##1{\textcolor[rgb]{0.73,0.13,0.13}{##1}}}
\@namedef{PY@tok@dl}{\def\PY@tc##1{\textcolor[rgb]{0.73,0.13,0.13}{##1}}}
\@namedef{PY@tok@s2}{\def\PY@tc##1{\textcolor[rgb]{0.73,0.13,0.13}{##1}}}
\@namedef{PY@tok@sh}{\def\PY@tc##1{\textcolor[rgb]{0.73,0.13,0.13}{##1}}}
\@namedef{PY@tok@s1}{\def\PY@tc##1{\textcolor[rgb]{0.73,0.13,0.13}{##1}}}
\@namedef{PY@tok@mb}{\def\PY@tc##1{\textcolor[rgb]{0.40,0.40,0.40}{##1}}}
\@namedef{PY@tok@mf}{\def\PY@tc##1{\textcolor[rgb]{0.40,0.40,0.40}{##1}}}
\@namedef{PY@tok@mh}{\def\PY@tc##1{\textcolor[rgb]{0.40,0.40,0.40}{##1}}}
\@namedef{PY@tok@mi}{\def\PY@tc##1{\textcolor[rgb]{0.40,0.40,0.40}{##1}}}
\@namedef{PY@tok@il}{\def\PY@tc##1{\textcolor[rgb]{0.40,0.40,0.40}{##1}}}
\@namedef{PY@tok@mo}{\def\PY@tc##1{\textcolor[rgb]{0.40,0.40,0.40}{##1}}}
\@namedef{PY@tok@ch}{\let\PY@it=\textit\def\PY@tc##1{\textcolor[rgb]{0.25,0.50,0.50}{##1}}}
\@namedef{PY@tok@cm}{\let\PY@it=\textit\def\PY@tc##1{\textcolor[rgb]{0.25,0.50,0.50}{##1}}}
\@namedef{PY@tok@cpf}{\let\PY@it=\textit\def\PY@tc##1{\textcolor[rgb]{0.25,0.50,0.50}{##1}}}
\@namedef{PY@tok@c1}{\let\PY@it=\textit\def\PY@tc##1{\textcolor[rgb]{0.25,0.50,0.50}{##1}}}
\@namedef{PY@tok@cs}{\let\PY@it=\textit\def\PY@tc##1{\textcolor[rgb]{0.25,0.50,0.50}{##1}}}

\def\PYZbs{\char`\\}
\def\PYZus{\char`\_}
\def\PYZob{\char`\{}
\def\PYZcb{\char`\}}
\def\PYZca{\char`\^}
\def\PYZam{\char`\&}
\def\PYZlt{\char`\<}
\def\PYZgt{\char`\>}
\def\PYZsh{\char`\#}
\def\PYZpc{\char`\%}
\def\PYZdl{\char`\$}
\def\PYZhy{\char`\-}
\def\PYZsq{\char`\'}
\def\PYZdq{\char`\"}
\def\PYZti{\char`\~}
% for compatibility with earlier versions
\def\PYZat{@}
\def\PYZlb{[}
\def\PYZrb{]}
\makeatother


    % For linebreaks inside Verbatim environment from package fancyvrb. 
    \makeatletter
        \newbox\Wrappedcontinuationbox 
        \newbox\Wrappedvisiblespacebox 
        \newcommand*\Wrappedvisiblespace {\textcolor{red}{\textvisiblespace}} 
        \newcommand*\Wrappedcontinuationsymbol {\textcolor{red}{\llap{\tiny$\m@th\hookrightarrow$}}} 
        \newcommand*\Wrappedcontinuationindent {3ex } 
        \newcommand*\Wrappedafterbreak {\kern\Wrappedcontinuationindent\copy\Wrappedcontinuationbox} 
        % Take advantage of the already applied Pygments mark-up to insert 
        % potential linebreaks for TeX processing. 
        %        {, <, #, %, $, ' and ": go to next line. 
        %        _, }, ^, &, >, - and ~: stay at end of broken line. 
        % Use of \textquotesingle for straight quote. 
        \newcommand*\Wrappedbreaksatspecials {% 
            \def\PYGZus{\discretionary{\char`\_}{\Wrappedafterbreak}{\char`\_}}% 
            \def\PYGZob{\discretionary{}{\Wrappedafterbreak\char`\{}{\char`\{}}% 
            \def\PYGZcb{\discretionary{\char`\}}{\Wrappedafterbreak}{\char`\}}}% 
            \def\PYGZca{\discretionary{\char`\^}{\Wrappedafterbreak}{\char`\^}}% 
            \def\PYGZam{\discretionary{\char`\&}{\Wrappedafterbreak}{\char`\&}}% 
            \def\PYGZlt{\discretionary{}{\Wrappedafterbreak\char`\<}{\char`\<}}% 
            \def\PYGZgt{\discretionary{\char`\>}{\Wrappedafterbreak}{\char`\>}}% 
            \def\PYGZsh{\discretionary{}{\Wrappedafterbreak\char`\#}{\char`\#}}% 
            \def\PYGZpc{\discretionary{}{\Wrappedafterbreak\char`\%}{\char`\%}}% 
            \def\PYGZdl{\discretionary{}{\Wrappedafterbreak\char`\$}{\char`\$}}% 
            \def\PYGZhy{\discretionary{\char`\-}{\Wrappedafterbreak}{\char`\-}}% 
            \def\PYGZsq{\discretionary{}{\Wrappedafterbreak\textquotesingle}{\textquotesingle}}% 
            \def\PYGZdq{\discretionary{}{\Wrappedafterbreak\char`\"}{\char`\"}}% 
            \def\PYGZti{\discretionary{\char`\~}{\Wrappedafterbreak}{\char`\~}}% 
        } 
        % Some characters . , ; ? ! / are not pygmentized. 
        % This macro makes them "active" and they will insert potential linebreaks 
        \newcommand*\Wrappedbreaksatpunct {% 
            \lccode`\~`\.\lowercase{\def~}{\discretionary{\hbox{\char`\.}}{\Wrappedafterbreak}{\hbox{\char`\.}}}% 
            \lccode`\~`\,\lowercase{\def~}{\discretionary{\hbox{\char`\,}}{\Wrappedafterbreak}{\hbox{\char`\,}}}% 
            \lccode`\~`\;\lowercase{\def~}{\discretionary{\hbox{\char`\;}}{\Wrappedafterbreak}{\hbox{\char`\;}}}% 
            \lccode`\~`\:\lowercase{\def~}{\discretionary{\hbox{\char`\:}}{\Wrappedafterbreak}{\hbox{\char`\:}}}% 
            \lccode`\~`\?\lowercase{\def~}{\discretionary{\hbox{\char`\?}}{\Wrappedafterbreak}{\hbox{\char`\?}}}% 
            \lccode`\~`\!\lowercase{\def~}{\discretionary{\hbox{\char`\!}}{\Wrappedafterbreak}{\hbox{\char`\!}}}% 
            \lccode`\~`\/\lowercase{\def~}{\discretionary{\hbox{\char`\/}}{\Wrappedafterbreak}{\hbox{\char`\/}}}% 
            \catcode`\.\active
            \catcode`\,\active 
            \catcode`\;\active
            \catcode`\:\active
            \catcode`\?\active
            \catcode`\!\active
            \catcode`\/\active 
            \lccode`\~`\~ 	
        }
    \makeatother

    \let\OriginalVerbatim=\Verbatim
    \makeatletter
    \renewcommand{\Verbatim}[1][1]{%
        %\parskip\z@skip
        \sbox\Wrappedcontinuationbox {\Wrappedcontinuationsymbol}%
        \sbox\Wrappedvisiblespacebox {\FV@SetupFont\Wrappedvisiblespace}%
        \def\FancyVerbFormatLine ##1{\hsize\linewidth
            \vtop{\raggedright\hyphenpenalty\z@\exhyphenpenalty\z@
                \doublehyphendemerits\z@\finalhyphendemerits\z@
                \strut ##1\strut}%
        }%
        % If the linebreak is at a space, the latter will be displayed as visible
        % space at end of first line, and a continuation symbol starts next line.
        % Stretch/shrink are however usually zero for typewriter font.
        \def\FV@Space {%
            \nobreak\hskip\z@ plus\fontdimen3\font minus\fontdimen4\font
            \discretionary{\copy\Wrappedvisiblespacebox}{\Wrappedafterbreak}
            {\kern\fontdimen2\font}%
        }%
        
        % Allow breaks at special characters using \PYG... macros.
        \Wrappedbreaksatspecials
        % Breaks at punctuation characters . , ; ? ! and / need catcode=\active 	
        \OriginalVerbatim[#1,codes*=\Wrappedbreaksatpunct]%
    }
    \makeatother

    % Exact colors from NB
    \definecolor{incolor}{HTML}{303F9F}
    \definecolor{outcolor}{HTML}{D84315}
    \definecolor{cellborder}{HTML}{CFCFCF}
    \definecolor{cellbackground}{HTML}{F7F7F7}
    
    % prompt
    \makeatletter
    \newcommand{\boxspacing}{\kern\kvtcb@left@rule\kern\kvtcb@boxsep}
    \makeatother
    \newcommand{\prompt}[4]{
        {\ttfamily\llap{{\color{#2}[#3]:\hspace{3pt}#4}}\vspace{-\baselineskip}}
    }
    

    
    % Prevent overflowing lines due to hard-to-break entities
    \sloppy 
    % Setup hyperref package
    \hypersetup{
      breaklinks=true,  % so long urls are correctly broken across lines
      colorlinks=true,
      urlcolor=urlcolor,
      linkcolor=linkcolor,
      citecolor=citecolor,
      }
    % Slightly bigger margins than the latex defaults
    
    \geometry{verbose,tmargin=1in,bmargin=1in,lmargin=1in,rmargin=1in}
    
    

\begin{document}
    
    \maketitle
    
    

    
    \section{Соответствуют ли уравнения для запаздывающих потенциалов
калибровке
Лоренца?}\label{ux441ux43eux43eux442ux432ux435ux442ux441ux442ux432ux443ux44eux442-ux43bux438-ux443ux440ux430ux432ux43dux435ux43dux438ux44f-ux434ux43bux44f-ux437ux430ux43fux430ux437ux434ux44bux432ux430ux44eux449ux438ux445-ux43fux43eux442ux435ux43dux446ux438ux430ux43bux43eux432-ux43aux430ux43bux438ux431ux440ux43eux432ux43aux435-ux43bux43eux440ux435ux43dux446ux430}

    А.Ю.Дроздов

    Поле векторного потенциала создаваемое движущимся равномерно и
прямолинейно зарядом, рассчитывается в классической электродинамике в
соответствии с
\[\overrightarrow{{{A}_{H}}}=\frac{1}{c}\int{\frac{\overrightarrow{j}\left( x',y',z',t-\frac{R}{v} \right)}{R}}\ dV'\]
Математически это выражение является решением уравнения Даламбера для
векторного потенциала. Представляет интерес задаться вопросом,
соответствует ли это выражение калибровке Лоренца? Для этого проследим
за выкладками {[} CITATION Там57 \l 1049 {]} §96, в котором он ищет
выражение для дивергенции векторного потенциала
\[{{div}_{a}}\ \overrightarrow{{{A}_{H}}}=\frac{1}{c}di{{v}_{a}}\ \int{\frac{\overrightarrow{j}\left( x',y',z',t-\frac{R}{v} \right)}{R}}\ dV'=\frac{1}{c}\ \int{di{{v}_{a}}\frac{\overrightarrow{j}\left( x',y',z',t-\frac{R}{v} \right)}{R}}\ dV'\]
Применяя уравнение векторного анализа
\[div\left( \varphi \overrightarrow{a} \right)=\varphi \ div\left( \overrightarrow{a} \right)+a\ grad\left( \varphi  \right)\]
\[{{div}_{a}}\frac{\overrightarrow{j}\left( x',y',z',t-\frac{R}{v} \right)}{R}=\frac{1}{R}\cdot di{{v}_{a}}\ \overrightarrow{j}\left( x',y',z',t-\frac{R}{v} \right)+\overrightarrow{j}\left( x',y',z',t-\frac{R}{v} \right)\cdot \nabla \frac{1}{R}\]
Так как аргумент вектора \(\overrightarrow{j}\) зависит от \(x,y,z\)
только через посредством эффективного времени \(t'=t-\frac{R}{v}\), то
\[{{div}_{a}}\ \overrightarrow{j}\left( x',y',z',t-\frac{R}{v} \right)=\frac{\partial {{j}_{x}}}{\partial t'}\frac{\partial t'}{\partial x}+\frac{\partial {{j}_{y}}}{\partial t'}\frac{\partial t'}{\partial y}+\frac{\partial {{j}_{z}}}{\partial t'}\frac{\partial t'}{\partial z}=\frac{\partial \overrightarrow{j}}{\partial t'}\cdot \nabla t'=\frac{\partial \overrightarrow{j}}{\partial t'}\cdot \nabla \left( t-\frac{R}{v} \right)=-\frac{1}{v}\frac{\partial \overrightarrow{j}}{\partial t'}\cdot \nabla R\]
Окончательно
\[{{div}_{a}}\frac{\overrightarrow{j}}{R}=\overrightarrow{j}\cdot \nabla \frac{1}{R}-\frac{1}{vR}\frac{\partial \overrightarrow{j}}{\partial t'}\cdot \nabla R\]
С другой стороны, очевидно, что
\(div{{'}_{q}}\ \frac{\overrightarrow{j}}{R}={{\left( div{{'}_{q}}\ \frac{\overrightarrow{j}}{R} \right)}_{t'=const}}+\frac{1}{R}\frac{\partial \overrightarrow{j}}{\partial t'}\cdot \nabla 't'\)
где \(div{{'}_{q}}\) и \(\nabla '\) в отличие от \({{div}_{a}}\) и
$\nabla$ означает дифференцирование по координатам истока
\(x',y',z'\). После преобразований
\[div{{'}_{q}}\frac{\overrightarrow{j}}{R}=\overrightarrow{j}\cdot \nabla '\frac{1}{R}+\frac{1}{R}{{\left( div{{'}_{q}}\ \overrightarrow{j} \right)}_{t'=const}}-\frac{1}{vR}\frac{\partial \overrightarrow{j}}{\partial t'}\cdot \nabla 'R\]
Приняв во внимание, что уравнение
\(gra{{d}_{a}}\ f\left( R \right)=-\ gra{{d}_{q}}\ f\left( R \right)\) в
теперешних обозначениях принимает вид
\[\nabla f\left( R \right)=-\ \nabla 'f\left( R \right)\] и сравнивая
выражения для \[di{{v}_{a}}\frac{\overrightarrow{j}}{R}\] и
\[div{{'}_{q}}\frac{\overrightarrow{j}}{R}\]
\[{{div}_{a}}\frac{\overrightarrow{j}}{R}=-\ div{{'}_{q}}\frac{\overrightarrow{j}}{R}+\frac{1}{R}{{\left( div{{'}_{q}}\ \overrightarrow{j} \right)}_{t'=const}}\]
Стало быть из 13 получено Таммом
\[{{div}_{a}}\ \overrightarrow{{{A}_{H}}}=-\frac{1}{c}\ \int{div{{'}_{q}}\frac{\overrightarrow{j}\left( x',y',z',t-\frac{R}{v} \right)}{R}}\ dV'+\frac{1}{c}\int{\frac{dV'}{R}}{{\left( div{{'}_{q}}\ \overrightarrow{j} \right)}_{t'=const}}\]
Первый из этих интегралов может быть преобразован с помощью теоремы
Гаусса в интеграл по поверхности \(S\), охватывающей объём \(V'\)
\[\int{div{{'}_{q}}\frac{\overrightarrow{j}\left( x',y',z',t-\frac{R}{v} \right)}{R}}\ dV'=\oint{\frac{{{j}_{n}}\left( x',y',z',t-\frac{R}{v} \right)}{R}}\ dS\]
Далее Тамм пишет: «если распространить интегрирование на все бесконечное
пространство, то этот интеграл обратится в нуль, если только все
электрические токи сосредоточены в конечной области пространства».
Однако аналогичный интеграл встречается у Тамма в §46, в котором
рассматривая пондемоторное взаимодействие постоянных токов, Тамм не
распространял интегрирование на все бесконечное пространство, а писал,
что «интегрирование должно быть распространено по поверхности всех
обтекаемых током проводников». И это верно в рамках представлений
классической электродинамики, поскольку классическая электродинамика при
вычислении векторного потенциала (как в калибровке Кулона, так и в
калибровке Лоренца) производит интегрирование только лишь по объёму
токов проводимости, занимающих конечный объём, но не производит
интегрирование по объёму токов смещения, занимающих все бесконечное
пространство. Следовательно, попытка Тамма при вычислении интеграла по
объёму токов проводимости «распространить интегрирование на все
бесконечное пространство» является, на мой взгляд, математически
некорректной. Самое интересное, что, перевернув страницу, в продолжении
доказательства Тамма можно встретить ещё одно утверждение, в
математической корректности которого есть определённые сомнения. Тамм
пишет: «обратимся теперь к правой части соотношения (96.3).
\[\varphi \left( x,y,z,t \right)=\frac{1}{\varepsilon }\int{\frac{\rho \left( x',y',z',t-\frac{R}{v} \right)}{R}}\ dV'\]
В виду уравнения (95.5) \(t'=t-\frac{R}{v}\) для произвольной функции
\(f\left( t' \right)\) имеем
\(\frac{\partial f\left( t' \right)}{\partial t}=\frac{\partial f\left( t' \right)}{\partial t'}\frac{\partial t'}{\partial t}=\frac{\partial f\left( t' \right)}{\partial t'}\)
Поэтому дифференцируя уравнение (96.3) по времени под знаком интеграла»,
Тамм получает:
\[-\frac{\varepsilon \mu }{c}\frac{\partial \varphi \left( x,y,z,t \right)}{\partial t}=-\frac{\mu }{c}\int{\frac{dV'}{R}\frac{\partial \rho \left( x',y',z',t-\frac{R}{v} \right)}{\partial t'}}\ \]
Здесь следует обратить внимание, что приравнивание
\(\frac{\partial t'}{\partial t}\) единице не корректно. Чтобы это
показать, найдём производную запаздывающего момента \(t'=t-\frac{R}{c}\)
по времени \[\frac{\partial t'}{\partial t}\] Дифференцируя
\[R\left( \overrightarrow{x},\ \overrightarrow{x'}\left( t' \right) \right)=c\left( t-t' \right)\]
по \(t\) получаем
\(\frac{\partial R}{\partial t}=\frac{\partial R}{\partial t'}\frac{\partial t'}{\partial t}=c\left( 1-\frac{\partial t'}{\partial t} \right)\).
Значение \(\frac{\partial R}{\partial t'}\) находим так
\[\frac{\partial R\left( t' \right)}{\partial t'}=\frac{\partial }{\partial t'}{{\left( \overrightarrow{R}\cdot \overrightarrow{R} \right)}^{{}^{1}/{}_{2}}}=\frac{1}{2}{{\left( \overrightarrow{R}\cdot \overrightarrow{R} \right)}^{-{}^{1}/{}_{2}}}\frac{\partial }{\partial t'}\left( \overrightarrow{R}\cdot \overrightarrow{R} \right)=\frac{1}{2R}\left( \overrightarrow{R}\cdot \frac{\partial \overrightarrow{R}}{\partial t'}+\frac{\partial \overrightarrow{R}}{\partial t'}\cdot \overrightarrow{R} \right)=\frac{\overrightarrow{R}\left( t' \right)}{R\left( t' \right)}\cdot \frac{\partial \overrightarrow{R}\left( t' \right)}{\partial t'}=-\frac{\overrightarrow{R}\cdot \overrightarrow{v}}{R}\]
Далее учитывая что
\(\frac{\partial \overrightarrow{R}}{\partial t'}=-\overrightarrow{v}\left( t' \right)\)
имеем
\[\frac{\partial R}{\partial t'}=-\frac{\overrightarrow{R}}{R}\cdot \overrightarrow{v}\]
теперь
\(\frac{\partial R}{\partial t'}\frac{\partial t'}{\partial t}=c\left( 1-\frac{\partial t'}{\partial t} \right)\)преобразуется
к
\(-\frac{\overrightarrow{R}\cdot \overrightarrow{v}}{cR}\frac{\partial t'}{\partial t}=1-\frac{\partial t'}{\partial t}\)
решая которое получаем
\[\frac{\partial t'}{\partial t}=\frac{1}{1-\frac{\overrightarrow{v}\cdot \overrightarrow{R}}{Rc}}\]
Отсюда правильно производную скалярного потенциала по времени записать
так
\[-\frac{\varepsilon \mu }{c}\frac{\partial \varphi \left( x,y,z,t \right)}{\partial t}=-\frac{\mu }{c}\int{\frac{dV'}{R-\frac{\overrightarrow{v}\cdot \overrightarrow{R}}{c}}\frac{\partial \rho \left( x',y',z',t-\frac{R}{v} \right)}{\partial t'}}\]
Применяя соотношение непрерывности в виде
\[{{\left( div{{'}_{q}}\ \overrightarrow{j} \right)}_{t'=const}}=-\frac{\partial \rho \left( x',y',z',t' \right)}{\partial t'}\],
убеждаемся что
\[\frac{1}{c}\ \int{div{{'}_{q}}\frac{\overrightarrow{j}\left( x',y',z',t-\frac{R}{v} \right)}{R}}\ dV'+di{{v}_{a}}\ \overrightarrow{{{A}_{H}}}=-\frac{1}{c}\int{\frac{dV'}{R}}\frac{\partial \rho \left( x',y',z',t' \right)}{\partial t'}\]
Допуская применимость калибровки Лоренца
\[div\ \overrightarrow{{{A}_{H}}}=-\frac{\varepsilon }{c}\frac{\partial \varphi }{\partial t}\]
запишем следующее соотношение
\[\frac{1}{c}\ \int{div{{'}_{q}}\frac{\overrightarrow{j}\left( x',y',z',t-\frac{R}{v} \right)}{R}}\ dV'-\frac{1}{c}\int{\frac{dV'}{R-\frac{\overrightarrow{v}\cdot \overrightarrow{R}}{c}}\frac{\partial \rho \left( x',y',z',t-\frac{R}{v} \right)}{\partial t'}}=-\frac{1}{c}\int{\frac{dV'}{R}}\frac{\partial \rho \left( x',y',z',t' \right)}{\partial t'}\]
Или
\[\frac{1}{c}\ \oint{\frac{{{j}_{n}}\left( x',y',z',t' \right)}{R}}\ dS'-\frac{1}{c}\int{\frac{dV'}{R-\frac{\overrightarrow{v}\cdot \overrightarrow{R}}{c}}\frac{\partial \rho \left( x',y',z',t' \right)}{\partial t'}}=-\frac{1}{c}\int{\frac{dV'}{R}}\frac{\partial \rho \left( x',y',z',t' \right)}{\partial t'}\]

Из вышеизложенного видно, что доказательство соответствия решения
уравнения Даламбера калибровке Лоренца, приведенное Таммом содержит две
математические некорректности. Из вышеизложенного становится очевидным,
что приём распространения интегрирования по токам проводимости на все
бесконечное пространство не верен. Потому что при обнулении интеграла по
\[dS'\] в последнем соотношении мы приходим к очевидному неравенству,
следствием которого становится неприменимость соотношения Лоренца для
потенциалов. С другой стороны, справедливость полученного в данной
работе соотношения 14 для произвольных токовых систем также не является
очевидной. Преобразуем этот интеграл
\[\ \oint{\frac{{{j}_{n}}\left( x',y',z',t' \right)}{R}}\ dS'=\int{\left( \frac{1}{R-\frac{\overrightarrow{v}\cdot \overrightarrow{R}}{c}}-\frac{1}{R} \right)\frac{\partial \rho \left( x',y',z',t' \right)}{\partial t'}}\ dV'=\int{\left( \frac{\overrightarrow{v}\cdot \overrightarrow{R}}{R\left( cR-\overrightarrow{v}\cdot \overrightarrow{R} \right)} \right)\frac{\partial \rho \left( x',y',z',t' \right)}{\partial t'}}\ dV'\]
Делаем обратную замену в левой части по тереме Гаусса
\[\int{div{{'}_{q}}\frac{\overrightarrow{j}\left( x',y',z',t' \right)}{R}}\ dV'=\int{\left( \frac{\overrightarrow{v}\cdot \overrightarrow{R}}{R\left( cR-\overrightarrow{v}\cdot \overrightarrow{R} \right)} \right)\frac{\partial \rho \left( x',y',z',t' \right)}{\partial t'}}\ dV'\]
Поскольку это равенство должно выполняться для любого объёма токовых
систем, то это даёт нам право приравнять подынтегральные выражения.
\[div{{'}_{q}}\frac{\overrightarrow{j}\left( x',y',z',t' \right)}{R}\ =\ \left( \frac{\overrightarrow{v}\cdot \overrightarrow{R}}{R\left( cR-\overrightarrow{v}\cdot \overrightarrow{R} \right)} \right)\frac{\partial \rho \left( x',y',z',t' \right)}{\partial t'}=\left( \frac{1}{R-\frac{\overrightarrow{v}\cdot \overrightarrow{R}}{c}}-\frac{1}{R} \right)\frac{\partial \rho \left( x',y',z',t' \right)}{\partial t'}\]
Физический смысл полученного выражения требует исследования. Левая часть
этого выражения уже преобразовывалась Таммом в следующем виде
\[div{{'}_{q}}\frac{\overrightarrow{j}}{R}=\overrightarrow{j}\cdot \nabla '\frac{1}{R}+\frac{1}{R}{{\left( div{{'}_{q}}\ \overrightarrow{j} \right)}_{t'=const}}-\frac{1}{vR}\frac{\partial \overrightarrow{j}}{\partial t'}\cdot \nabla 'R\]
Опять применяя соотношение непрерывности в виде
\[{{\left( div{{'}_{q}}\ \overrightarrow{j} \right)}_{t'=const}}=-\frac{\partial \rho \left( x',y',z',t' \right)}{\partial t'}\],
и возвращая правую часть равенства получаем
\[\overrightarrow{j}\cdot \nabla '\frac{1}{R}-\frac{1}{R}\frac{\partial \rho \left( x',y',z',t' \right)}{\partial t'}-\frac{1}{vR}\frac{\partial \overrightarrow{j}}{\partial t'}\cdot \nabla 'R=\left( \frac{1}{R-\frac{\overrightarrow{v}\cdot \overrightarrow{R}}{c}}-\frac{1}{R} \right)\frac{\partial \rho \left( x',y',z',t' \right)}{\partial t'}\]
Сокращаем
\[\overrightarrow{j}\cdot \nabla '\frac{1}{R}-\frac{1}{vR}\frac{\partial \overrightarrow{j}}{\partial t'}\cdot \nabla 'R=\left( \frac{1}{R-\frac{\overrightarrow{v}\cdot \overrightarrow{R}}{c}} \right)\frac{\partial \rho \left( x',y',z',t' \right)}{\partial t'}\]

    \begin{tcolorbox}[breakable, size=fbox, boxrule=1pt, pad at break*=1mm,colback=cellbackground, colframe=cellborder]
\prompt{In}{incolor}{ }{\boxspacing}
\begin{Verbatim}[commandchars=\\\{\}]

\end{Verbatim}
\end{tcolorbox}


    % Add a bibliography block to the postdoc
    
    
    
\end{document}
